\documentclass{article}
\usepackage[a4paper, total={6in, 8in}]{geometry}
\usepackage{graphicx} % Required for inserting images
\usepackage{amsmath,amssymb,xcolor,bm,amsthm,mathtools,comment, enumerate}
\newcommand{\eps}{\varepsilon}
\usepackage{tikz}
\newtheorem{theorem}{Theorem}[section]
\newtheorem{proposition}[theorem]{Proposition}
\newtheorem{remark}[theorem]{Remark}
\newtheorem{definition}[theorem]{Definition}
\newtheorem{lemma}[theorem]{Lemma}
\newtheorem{corollary}[theorem]{Corollary}
\newtheorem{example}[theorem]{Example}
\newtheorem{assumption}[theorem]{Assumption}
\newtheorem{conjecture}[theorem]{Conjecture}
\newtheorem{question}{Question}

\title{Every tuple of natural numbers of norm below 2 to the n minus one is an admissible profile}
\author{Joshua Erde}
\date{arxiv:2404.03950v3; CC BY 4.0}
\begin{document}
\maketitle
\noindent\textit{Editorial note.} The counted unit is the article's own Theorem c:admissible at source lines 292-294 together with the proof environment that begins on the next line (295-300); the pairing is the ordinary adjacent one.\par\medskip
\noindent\textit{Editorial note.} The retained context at source lines 30-36 and 51 is the article's own statement of the Balister-Gyori-Schelp necessary condition and conjecture and its own definitions of profile, admissible, even and perfect tuples.\par\medskip
\noindent\textit{Editorial note.} The frozen copy the spans of this card are taken from is the fl0X staging of the same source\_id, that is sources/combinatorics-fl03-matchings-in-the-hypercube-with-sp, whose licence observation carries status supported; the shell directory named after the arXiv id holds no licence evidence of its own.\par\medskip
\noindent\textit{Editorial note.} The two bibliography entries retained with this card, HKAES23 and HJR52, are copied verbatim from the author's own compiled bibliography main.bbl as shipped in the pinned e-print of arXiv:2404.03950v3; no entry was reworded, added or reordered, and the file main.tex of the same e-print shows the source cites exactly these two keys in this proof.\par\medskip
\noindent\textit{Editorial note.} The author's own acknowledgements at source line 353 record that this proof was provided to him by the anonymous referees; the counted unit is nonetheless the author's published text, retained line for line.\par\medskip
\noindent\textit{Editorial note.} The article's Theorem t:perfect at source lines 53-55, its Lemma l:admissible at 67-69 and its Lemma l:first at 154-156 are not counted here and are not used as premises.\par\medskip
\noindent\textit{Allowed standard prerequisites (editorial).}\begin{itemize}
\item Profiles of matchings in the hypercube Q\textasciicircum{}n and admissibility. For a matching M in the cube Q\textasciicircum{}n the profile is the vector x in N\textasciicircum{}n whose i-th coordinate counts the edges of M whose endpoints differ in the i-th coordinate; a tuple is admissible when it is the profile of some matching, a tuple is even when every coordinate is even, and it is perfect when the sum of its coordinates equals 2\textasciicircum{}\{n-1\}; the article's own definitions are retained at source line 51.. Fixes the meaning of admissible, which is exactly the conclusion of the counted theorem for every tuple of norm below 2\textasciicircum{}\{n-1\}.
\item The Balister-Gyori-Schelp partitioning problem and its vanishing-sum condition. Given elements d\_1, ..., d\_\{2\textasciicircum{}\{n-1\}\} of the vector space over F\_2 of dimension n, if the cube is partitioned into pairs with those differences then the sum of all d\_i must vanish; the article records this identity at source lines 30-33 and the conjecture that the condition is also sufficient at source lines 34-36.. The proof of the counted theorem invokes the conjecture for sequences all of whose differences have Hamming weight one, so both the statement and the surrounding setting must be available.
\item Hall's theorem on abelian groups and the batch-code argument of Hollmann, Khathuria, Riet and Skachek. The proof uses that a result of Hall implies the Balister-Gyori-Schelp conjecture for sequences in which every difference has odd Hamming weight, as recorded by Hollmann, Khathuria, Riet and Skachek in the remark cited by the author; the two bibliography entries are retained verbatim with this card.. This is the external fact that the author invokes at the opening of his proof; the counted unit is the derivation of the norm-below-threshold statement from it, not the fact itself.
\end{itemize}\medskip
\begin{equation}\label{e:null}
\sum_{i=1}^{2^{n-1}} \bm{d_i} = \sum_{i=1}^{2^{n-1}}  \bm{a_i} - \bm{b_i} =   \sum_{i=1}^{2^{n-1}}  \bm{a_i} + \bm{b_i} = \sum_{\bm{v} \in \mathbb{F}_2^n} \bm{v} = \bm{0}.
\end{equation}
They conjectured that this necessary condition was also sufficient.
\begin{conjecture}[Balister, Gy{\H{o}}ri, and Schelp]\label{c:BGS}
Let $n \geq 2$ and let $\bm{d_1},\ldots, \bm{d_{2^{n-1}}}$ be elements of $\mathbb{F}_2^n$ such that $\sum_{i=1}^{2^{n-1}} \bm{d_i} = \bm{0}$. Then there exists a partition of $\mathbb{F}_2^n$ into pairs $\{\bm{a_i},\bm{b_i}\}$ such that $\bm{a_i} - \bm{b_i} = \bm{d_i}$ for each $i \in \left[2^{n-1}\right]$.
\end{conjecture}

Given a matching $M \subseteq Q^n$ we say the \emph{profile} of $M$ is the vector $\bm{x}=(x_1,\ldots, x_n) \in \mathbb{N}^n$ such that $M$ contains $x_i$ edges whose endpoints differ in the $i$th coordinate, and we are interested in which tuples $\bm{x} \in \mathbb{N}^n$ are achievable as the profile of some matching. We call such tuples \emph{admissible}. In this case, the obvious necessary condition for a tuple to be the profile of a perfect matching is that each unit vector appears an even number of times as a matching edge, i.e., each $x_i$ is even, in which case we say that $\bm{x}$ is \emph{even}, and Conjecture \ref{c:BGS} would imply that this condition is also sufficient. The aim of this short note is to verify the conjecture in this case. Let us call a tuple $\bm{x} \in \mathbb{N}^n$  \emph{perfect} if $||\bm{x}||_1 := \sum_i x_i = 2^{n-1}$.

\begin{theorem}\label{c:admissible}
Let $n \in \mathbb{N}$. Every tuple $\bm{x} \in \mathbb{N}^n$ with $||\bm{x}||_1 < 2^{n-1}$ is admissible.
\end{theorem}

\begin{proof}
We first note that, by similar arguments as in \cite[Remark 5.4]{HKAES23}, a result of Hall \cite{HJR52} implies that Conjecture \ref{c:BGS} holds if all the vectors $\bm{d}_1,\ldots,\bm{d}^{2^{n-1}}$ have odd Hamming weight. 

Clearly it suffices to prove the statement for $\bm{x} \in \mathbb{N}^n$ with $||\bm{x}||_1 = 2^{n-1}-1$, and given any such vector $\bm{x}$ we can consider the sequence $\bm{d}_1,\ldots,\bm{d}^{2^{n-1}}$ given by $x_j$ copies of $\bm{e}_j$ for each $j \in [n]$ together with $\bm{d}_{2^{n-1}} = \sum_{i=1}^{2^{n-1}-1} \bm{d}_i$. Since each $\bm{d}_i$ with $i \in \left[2^{n-1}-1\right]$ has Hamming weight one, it follows that each $\bm{d}_i$ with $i\in \left[2^{n-1}\right]$ has odd Hamming weight, and hence there is a partition of $\mathbb{F}^n_2$ into pairs $\{\bm{a}_i,\bm{b}_i\}$ such that $\bm{a}_i - \bm{b}_i = \bm{d}_i$ for all $i \in \left[2^{n-1}\right]$. The matching $
M = \left\{ (\bm{a}_i,\bm{b}_i) \colon  \in \left[2^{n-1}-1\right]\right\}$ then shows that $\bm{x}$ is admissible.
\end{proof}

\begin{thebibliography}{99}
\bibitem{HJR52}
M.~Jr. Hall.
\newblock A combinatorial problem on {Abelian} groups.
\newblock {\em Proc. Am. Math. Soc.}, 3:584--587, 1952.

\bibitem{HKAES23}
H.~D.~L. Hollmann, K.~Khathuria, A.~Riet, and V.~Skachek.
\newblock On some batch code properties of the simplex code.
\newblock {\em Des. Codes Cryptography}, 91(5):1595--1605, 2023.
\end{thebibliography}
\end{document}
